\section{A direct quartic in 478 integer variables}
\label{sec:direct478}

The preceding compact equations optimize the number of integer unknowns.
There is also a concrete improvement at degree four using the separate
direct C222 proof-certificate source.  Let $R_0,\ldots,R_{304}$ be the
literal integral quadratic polynomials in
\texttt{direct/equation.json}, in its specified coordinate order, and set
\begin{equation}\label{eq:direct478}
 \boxed{P_{478}(z_0,\ldots,z_{477})=\sum_{j=0}^{304}R_j(z)^2.}
\end{equation}
The companion \texttt{expanded\_polynomial.json} contains every collected
coefficient.  A term $[c,[i_1,\ldots,i_d]]$ denotes
$cz_{i_1}\cdots z_{i_d}$, with an empty index list denoting the constant
monomial.  Every coefficient is a literal integer, and every variable
ranges independently over $\Z$.

\begin{theorem}\label{thm:direct478}
The polynomial \eqref{eq:direct478} has exact degree four, 478 occurring
variables, 5,484 collected monomials and coefficient height 511,326.
The ordinary proofs below, composed with the explicitly retained direct
C222/refutation calibration, establish
\[
 \mathrm{ZFC}\vdash
 \bigl[\neg\exists z\in\Z^{478}\ P_{478}(z)=0\bigr]
 \ \longleftrightarrow\ \operatorname{Con}(\mathrm{ZFC}).
\]
\end{theorem}

The matched predecessor is the derived-label square-shell endpoint in the
incoming powered-saturation/factorial-carrier archive at the comparison
snapshot.  Its metrics, and the new metrics, are:
\begin{center}
\begin{tabular}{lrrrr}
 & Variables & Quadratic rows & Monomials & Height\\
\hline
Matched predecessor & 517 & 310 & 6,774 & 3,904,267\\
New direct equation & 478 & 305 & 5,484 & 511,326
\end{tabular}
\end{center}
This improves all four displayed metrics for that direct branch.  Other
repository constructions have fewer variables with different degree or
coefficient costs.  No global minimum is asserted.

\subsection{The dependency order behind the saving}

The source supplies a proposed factorial $F$, a horizon $t=t_0+16$,
a proposed power $V$, and exact polynomial relations
\begin{align}
 &D=r+s+1, & DF&=V+s,\\
 &q=D+o_0+1, &q&=B+o_1+1,\\
 &B i=C+F^2+1, &&C\ge1.
\end{align}
Here $D$ and $B$ are proposed binomial outputs.  Their eventual meanings
are not assumed in these identities.  The power interface requests
\[
 V=(Ft^2+1)^t,\qquad N=B^K,\qquad q=N(N+V+1)^K,
 \qquad K=2(2+t+x+y+F),
\]
as well as the powers used for digit and range extraction.
Before power recovery, retain $t_0,x\ge0$, the endpoint and encoded-input
signs, and $o_0\ge0$.  Strengthen the existing factorial guard to $F\ge1$.
It costs the same number of witnesses and holds for every genuine factorial.

\begin{lemma}[Factorial signs after one power]\label{lem:early-fact-signs}
Suppose $F\ge1$, $V>F$, $D=r+s+1$, $DF=V+s$ and $rs\ge0$.
Then $r,s\ge0$ and $D\ge1$.
\end{lemma}
\begin{proof}
Eliminate $D$ to obtain $V=Fr+(F-1)s+F$.  If either member of $(r,s)$
were negative, product nonnegativity would make both nonpositive.
The displayed expression would then give $V\le F$, a contradiction.
This includes the cases where one member is zero.
\end{proof}

The base $Ft^2+1$ is already at least two.  The signed power certificate,
whose frame is justified below independently of these signs, therefore
recovers $V=(Ft^2+1)^t>F$ first.  Lemma~\ref{lem:early-fact-signs} proves
$D\ge1$, and hence $q=D+o_0+1\ge2$.  Next require only $Bo_1\ge0$.
The sum $B+o_1=q-1>0$ forces $B,o_1\ge0$; the reciprocal index identity
excludes $B=0$, so $B\ge1$ before its power is interpreted.

The two calls at index $K$ now orient their shared Pell pair.  The first
either gives $N=B^K$ or has the exceptional inverse case $B=1,N=-1$.
In both cases its second base $N+V+1$ is at least two.  That second call
fixes the orientation and excludes the inverse case.  All remaining seed
pairs have supporting calls at $q$ or $q+1$, both at least two.  The
unchanged addition formula recovers the single derived $t+1$ pair.
Thus all nineteen requested powers have their intended values.

Only now recover the other delayed signs.  Both binomial remainder pairs
and all ten range remainder pairs have exact sum equations
\[
 r_j+s_j=q^{\ell_j}-1>0,
\]
where the already positive exponent is $t,t+1,x+1$, or $y+1$.
For each such pair, $r_js_j\ge0$ implies both coordinates are nonnegative.
This replaces fourteen pairs of scalar predicates by fourteen product
predicates.  No binomial, factorial or CRT interpretation is used to
restore a sign that its own interpretation needs.

For unrestricted integer witnesses, use
\[
 Q(a,b,c)=a^2+b^2+c(c+1).
\]
It represents exactly the nonnegative integers.  Nonnegativity is
immediate, and surjectivity follows from the three-square theorem applied
to $4n+1$: exactly one of the three squares is odd.
Consequently $rs=Q(a,b,c)$ is one quadratic row expressing $rs\ge0$.
The ordinary positive-sum/product observation is elementary and already
used in earlier repository work; the new saving is its dependency-ordered
application to this stronger source, enabled by the earlier factorial
bootstrap.

\subsection{One guard for orientation and a nonzero shared ordinate}

Let $A$ be the main Pell base, and write the shared auxiliary ordinate as
\[
 V_* = \left(\prod_iY_i^2\right)(v+1).
\]
The retained seed slacks give $Y_i\ge k_i\ge1$ before any power is
interpreted.  Separate guards on the frame orientation and on $v$ can be
replaced by two quadratic rows involving four fresh coordinates:
\begin{equation}\label{eq:joint-anchor}
 u=(v+1)^2,\qquad
 (A-1)u-1=\alpha^2+\beta^2+\gamma(\gamma+1).
\end{equation}
Since the right side is nonnegative, $(A-1)u\ge1$.  Thus $A\ge2$,
$u\ge1$ and $v\ne-1$, so $V_*\ne0$.  The sign of $V_*$ is immaterial.
Conversely, canonical positive Pell witnesses have $A\ge2,v\ge0$;
the left side in the second equation is nonnegative and admits the required
three-square witnesses.  This joins two useful conclusions in one predicate
while keeping the final equation quartic.

The unchanged auxiliary norm gives
\[
 U^2-(A^2-1)V_*^2=1,\qquad |U|\ge A.
\]
The comparison base $b$ satisfies
\[
 b=1+4V_*(q_4+1),\qquad b\equiv A\pmod U.
\]
We permit $q_4$ to have either sign.  The first equation gives
$b\equiv1\pmod4$; the second excludes $b=1$ because
$0<A-1<|U|$.  Therefore $|b|\ge3$.

The signed comparison-base index theorem is inherited from the incoming
quadratic-frames package, and its complete proofs are retained under
\texttt{direct/source/}.  For completeness, the mechanism is as follows.
Classify absolute Pell solutions at $A$ and signed solutions at $b$;
polynomial dependence in the base transfers their congruences.
Squared divisibility $Y_i^2\mid V_*$ gives $Y_i\mid h$ for the common
auxiliary index $h$.  The abscissa residue identity gives the comparison
index congruent to either sign of the main index modulo $2h$.
The base-one ordinate congruence gives either sign of $k_i$ modulo $4Y_i$.
Hence the main index is congruent to $\pm k_i$ modulo $2Y_i$.
Both lie in $[1,Y_i]$, so they agree; the possible common endpoint also
gives equality.  Thus
\[
 |X_i|=\chi_A(k_i),\qquad Y_i=\psi_A(k_i).
\]
This argument allows both signs of $V_*,U,b$ and the comparison coordinates.
It never requires the removed private signs to be pointwise positive.

\subsection{A frame that uses proved signs efficiently}

The original power frame squared every atom.  The new frame uses
\begin{equation}\label{eq:typed-frame}
 W=2+\sum_{a\text{ already proved }\ge0}a
       +\sum_{a\text{ of undetermined sign}}a^2.
\end{equation}
There are 41 distinct affine atoms: 18 enter linearly and 23 quadratically.
Their signs are certified by 163 ordered polynomial deductions from 31
retained scalar predicates.  The frame-definition row and every power
interpretation are excluded from those deductions.
For every integer $a$, $a^2\ge|a|$; hence $W\ge2$ and $W>|a|$ for every
base, exponent and proposed output in the finite power interface.

Writing $Z=Z_{\rm minus1}+1$, the unchanged frame norm is
\[
 A^2-W^3(W+2)Z^2=1.
\]
Use $\chi_a(j)+\psi_a(j)\sqrt{a^2-1}=(a+\sqrt{a^2-1})^j$.
Equation~\eqref{eq:joint-anchor} gives $A\ge2$, so $Z\ne0$ even though
its sign is unrestricted.  Pell classification at base $W+1$ and the
recurrence $\psi_{W+1}(h)\equiv h\pmod W$ yield
$A=\chi_{W+1}(h)$ with $h\ge W$.  Thus
\[
 A\ge(W+1)^W,\qquad
 T_n=2An-n^2-1>W^W+W \quad(1\le n<W).
\]
For example $T_n\ge2A-2$, and the binomial expansion gives
$(W+1)^W\ge2W^W$.  These bounds hold before any proposed output is
interpreted or given a sign.

At an indexed pair, the power row
\[
 X+(n-A)Y=m+T_nw
\]
gives $m\equiv n^k\pmod{T_n}$ for positive $X$, and
$mn^k\equiv-1\pmod{T_n}$ for negative $X$.  The frame bounds give
$|m-n^k|<T_n$ and $|mn^k+1|<T_n$.  A base $n\ge2$ excludes the
inverse case and gives $m=n^k$ exactly, with no prior sign assumption
on $m$.  This licenses the initial factorial call and the dependency order
used above.

\subsection{Both directions and the finite equation}

Given a zero of the new system, the preceding arguments recover all powers
and all semantic front signs.  The inherited self-bounding factorial,
exact digit, reciprocal prime-power CRT, derived-label bound and
pairwise-coprime Horn-flow proofs then recover a finite C222 refutation.
Their full ordinary proof texts are included under \texttt{direct/source/}.
The endpoint composition is
\texttt{KEndpoint.c222\_iff\_zf\_refutation} followed by
\texttt{ConstructibleChoice.zf\_refutation\_iff\_zfc\_refutation}.

Conversely, start with the inherited arithmetic front supplied by any
finite refutation.  It has $F=t!\ge1$ and nonnegative entries in every
fused pair.  Construct its new $W$ by \eqref{eq:typed-frame}, choose the
positive Pell solution at index $W$ for base $W+1$, and choose the main
pairs at their required indices.  The order of a unit in a finite
quadratic ring supplies a common positive auxiliary iterate divisible by
the entire product $\prod_iY_i^2$.  CRT supplies a positive comparison
base with the required two congruences.  Complete all comparison pairs
and power quotients.  These canonical witnesses have $v\ge0$, so
\eqref{eq:joint-anchor} and every product predicate admit their ordinary
three-square witnesses.  No consistency assumption is used in this map.

These are equivalences of existential zero assertions.  New private
comparison coordinates need not solve the old positive-domain auxiliary
rows; fresh old witnesses can be reconstructed after recovering the front.

The complete new source has 690 integer coordinates and 515 quadratic
rows.  The projection certificate records 210 integral monic restorations;
two remaining unused coordinates are discarded with restoration zero.
Every substitution keeps every residual quadratic.  This leaves 478
coordinates and 305 rows.  The final explicitly recorded integral
translations are bijections.  Expanding the sum of squares gives all 5,484
nonzero terms of \eqref{eq:direct478} and proves the stated exact metrics.

The portable builder regenerates all ten normative data files byte for
byte.  One independent checker reconstructs every guard and frame row;
a second checks all source hashes, transformations, projection inverses,
translations and collected quartic coefficients.  Both pass under normal
and optimized Python.  Their finite boundary controls support transcription
and degeneration checks; they neither decide solvability nor replace the
unbounded mathematical proof.  This branch uses the direct calculus,
without a K215 execution encoding or Jones universal-polynomial parameter.
